\section{Ethics and Privacy}\label{sec:ethics}

\paragraph{Bystanders.} A head-worn camera sees people who did not consent to being recorded.
\begin{itemize}
\item SpatialNet runs on the glasses, so localisation needs no frame to leave the device.
\item A downstream step can be limited to the target crop and to moments when the wearer points.
\item SpatialNet is trained to ignore other people's hands. It does not detect, identify or track people.
\end{itemize}
On-device processing still does not remove the need for visible recording indicators; the developer unit has user- and bystander-facing LEDs \tbd{confirm whether this may be stated}.

\paragraph{Data.} EgoPoint is used under its licence \tbd{name the licence; state whether the converted labels may be shared}. The images shown in this paper come from its test and training splits. \tbd{T5 and T6: describe consent, the ethics approval or exemption (with reference number after review), data retention, and the avoidance of bystanders' faces.}

\paragraph{Availability.} \tbd{Anonymised link to the code, conversion script, trained weights and evaluation outputs.}

\section{Conclusion}\label{sec:conclusion}
SpatialNet localises the object a smart-glasses wearer points at, from a single frame, with 0.72M parameters and one forward pass. It does so by chaining instance-anchored heads (hand, fingertip, direction, target) and by conditioning the target search on ray coordinates derived from its own predictions. Training with pasted, unlabelled hands teaches it to follow the wearer's hand. \tbd{One sentence with the key T1 and T2 numbers.} The main open problems are pointing scenes that EgoPoint does not contain, real bystander hands, and far targets along ambiguous rays. We expect progress on them to come from richer egocentric pointing data at least as much as from larger models.
